\begin{ex}
Đoạn thẳng AB dài 16 cm, M là trung điểm của AB. Độ dài đoạn AM là bao nhiêu xăng-ti-mét?
\shortans{8}
\loigiai{
$16 : 2 = 8$ cm.
}
\end{ex}

\begin{ex}
Một hình tròn có bán kính 9 cm. Đường kính hình tròn đó là bao nhiêu xăng-ti-mét?
\shortans{18}
\loigiai{
$9 \times 2 = 18$ cm.
}
\end{ex}

\begin{ex}
Một hình tròn có đường kính 30 cm. Bán kính hình tròn đó là bao nhiêu xăng-ti-mét?
\shortans{15}
\loigiai{
$30 : 2 = 15$ cm.
}
\end{ex}

\begin{ex}
Hình vuông có cạnh dài 6 cm. Chu vi của hình vuông đó là bao nhiêu xăng-ti-mét?
\shortans{24}
\loigiai{
$6 \times 4 = 24$ cm.
}
\end{ex}

\begin{ex}
Hình chữ nhật có chiều dài 12 cm, chiều rộng 8 cm. Chu vi của hình chữ nhật đó là bao nhiêu xăng-ti-mét?
\shortans{40}
\loigiai{
$(12 + 8) \times 2 = 40$ cm.
}
\end{ex}

\begin{ex}
Một hình tam giác có độ dài 3 cạnh là 9 cm, 12 cm, 15 cm. Chu vi tam giác đó là bao nhiêu xăng-ti-mét?
\shortans{36}
\loigiai{
$9 + 12 + 15 = 36$ cm.
}
\end{ex}

\begin{ex}
Một khối lập phương có bao nhiêu mặt?
\shortans{6}
\loigiai{
Có 6 mặt.
}
\end{ex}

\begin{ex}
Một khối lập phương có bao nhiêu cạnh?
\shortans{12}
\loigiai{
Có 12 cạnh.
}
\end{ex}

\begin{ex}
Một khối hộp chữ nhật có bao nhiêu đỉnh?
\shortans{8}
\loigiai{
Có 8 đỉnh.
}
\end{ex}

\begin{ex}
Đổi: $7\text{ cm } 4\text{ mm} = \dots\text{ mm}$.
\shortans{74}
\loigiai{
$70 + 4 = 74$ mm.
}
\end{ex}

\begin{ex}
Đổi: $1\text{ kg } = \dots\text{ g}$.
\shortans{1000}
\loigiai{
$1\text{ kg} = 1000\text{ g}$.
}
\end{ex}

\begin{ex}
Đổi: $1\text{ l } = \dots\text{ ml}$.
\shortans{1000}
\loigiai{
$1\text{ l} = 1000\text{ ml}$.
}
\end{ex}

\begin{ex}
Một chai sữa chứa 250 ml, 4 chai như thế chứa bao nhiêu mi-li-lít sữa?
\shortans{1000}
\loigiai{
$250 \times 4 = 1000$ ml.
}
\end{ex}

\begin{ex}
Một quả dưa nặng 1 kg 400 g. Quả dưa đó nặng bao nhiêu gam?
\shortans{1400}
\loigiai{
$1000 + 400 = 1400$ g.
}
\end{ex}

\begin{ex}
Nhiệt độ phòng là $24^\circ$C, sau đó tăng thêm $5^\circ$C. Nhiệt độ lúc sau là bao nhiêu độ C?
\shortans{29}
\loigiai{
$24 + 5 = 29^\circ$C.
}
\end{ex}

\begin{ex}
Một hình vuông có chu vi 32 cm. Độ dài mỗi cạnh của hình vuông là bao nhiêu xăng-ti-mét?
\shortans{8}
\loigiai{
$32 : 4 = 8$ cm.
}
\end{ex}

\begin{ex}
Một hình chữ nhật có chu vi 50 cm, chiều dài 15 cm. Chiều rộng hình chữ nhật đó là bao nhiêu xăng-ti-mét?
\shortans{10}
\loigiai{
Nửa chu vi = 25 cm; chiều rộng = $25 - 15 = 10$ cm.
}
\end{ex}

\begin{ex}
Dùng ê-ke để kiểm tra, hình chữ nhật có bao nhiêu góc vuông?
\shortans{4}
\loigiai{
Có 4 góc vuông.
}
\end{ex}

\begin{ex}
Một khối lập phương có cạnh dài 4 cm. Tổng độ dài tất cả 12 cạnh là bao nhiêu xăng-ti-mét?
\shortans{48}
\loigiai{
$4 \times 12 = 48$ cm.
}
\end{ex}

\begin{ex}
Đổi: $3\text{ dm } 5\text{ mm} = \dots\text{ mm}$.
\shortans{305}
\loigiai{
$300 + 5 = 305$ mm.
}
\end{ex}

\begin{ex}
Bác thợ rèn thanh sắt dài 1 m thành một khung hình vuông. Mỗi cạnh của khung hình vuông dài bao nhiêu xăng-ti-mét?
\shortans{25}
\loigiai{
$1\text{ m} = 100\text{ cm}$; $100 : 4 = 25$ cm.
}
\end{ex}

\begin{ex}
Một gói mì chính nặng 450 g, mẹ đã dùng 150 g. Còn lại bao nhiêu gam?
\shortans{300}
\loigiai{
$450 - 150 = 300$ g.
}
\end{ex}

\begin{ex}
Nhiệt kế chỉ $36^\circ$C. Vào ban đêm nhiệt độ giảm $7^\circ$C. Ban đêm nhiệt độ là bao nhiêu độ C?
\shortans{29}
\loigiai{
$36 - 7 = 29^\circ$C.
}
\end{ex}

\begin{ex}
Tính chu vi hình tam giác đều có cạnh 11 cm theo xăng-ti-mét.
\shortans{33}
\loigiai{
$11 \times 3 = 33$ cm.
}
\end{ex}

\begin{ex}
Một hình tứ giác có 4 cạnh đều bằng 7 cm. Chu vi của tứ giác đó là bao nhiêu xăng-ti-mét?
\shortans{28}
\loigiai{
$7 \times 4 = 28$ cm.
}
\end{ex}
